\section*{第 \ref{ch:b2:plant-flowering} 章——生殖发育与开花调控}

\begin{solution}{exo:b2:plant-flowering:1}
营养\omterm{def:b2:plant-meristems:meristem}{分生组织}：不断形成带有腋芽的叶，
为无限生长。\omterm{def:b2:plant-flowering:transition}{花序分生组织}：形成苞片和花\omterm{def:b2:plant-meristems:meristem}{分生组织}，
二者位于苞片腋内，多数物种中为无限生长。花\omterm{def:b2:plant-meristems:meristem}{分生组织}：
形成四轮花器官后停止生长，属有限生长。
\end{solution}

\begin{solution}{exo:b2:plant-flowering:2}
\omterm{prop:b2:plant-flowering:photoperiod}{短日植物}在日长短于临界长度时开花（菊花、大豆、水稻）；\omterm{prop:b2:plant-flowering:photoperiod}{长日植物}
在日长长于临界长度时开花（菠菜、小麦、
\emph{Arabidopsis}）；日中性植物
在长到一定大小时开花（番茄、玉米）。\omterm{prop:b2:plant-flowering:photoperiod}{短日植物}
测量的是不间断黑夜的长度。
\end{solution}

\begin{solution}{exo:b2:plant-flowering:3}
野生型：萼片、花瓣、\omterm{def:b2:angiosperm-reproduction:flower}{雄蕊}、\omterm{def:b2:angiosperm-reproduction:flower}{心皮}。无 A：\omterm{def:b2:angiosperm-reproduction:flower}{心皮}、\omterm{def:b2:angiosperm-reproduction:flower}{雄蕊}、
\omterm{def:b2:angiosperm-reproduction:flower}{雄蕊}、\omterm{def:b2:angiosperm-reproduction:flower}{心皮}。无 B：萼片、萼片、\omterm{def:b2:angiosperm-reproduction:flower}{心皮}、\omterm{def:b2:angiosperm-reproduction:flower}{心皮}。无 C：萼片、
花瓣、花瓣、萼片，且花在内部重复出现。
\end{solution}

\begin{solution}{exo:b2:plant-flowering:4}
经历数周的低温潮湿（冬季已过）；光照（\omterm{def:b2:angiosperm-reproduction:seed}{种子}位于
近地表处）；温度交替变化（裸露土壤，无植被遮盖）；
雨水浸出（雨季已至）；磨损处理（经过
动物肠道或土壤中的磨蚀）；烟雾与高温（火灾
已清理地面）。
\end{solution}

\begin{solution}{exo:b2:plant-flowering:5}
16/8：黑夜 8 小时，低于临界值 10 小时：不开花。12/12：开花。
12 小时光照、6 小时黑暗、闪光、6 小时黑暗：两段 6 小时的黑夜：不开花。
8 小时光照、4 小时黑暗、4 小时光照、8 小时黑暗：最长黑夜为 8 小时：不开花。
\end{solution}

\begin{solution}{exo:b2:plant-flowering:6}
\omterm{def:b2:plant-flowering:florigen}{成花素}是一种可通过嫁接传递的信号，由叶片产生，
且在\omterm{prop:b2:plant-flowering:photoperiod}{短日植物}和\omterm{prop:b2:plant-flowering:photoperiod}{长日植物}中相同——差别在于叶片何时产生它，
而不在于产生什么。切断韧皮部后，嫁接
便不能诱导开花：信号通过
韧皮部运输。
\end{solution}

\begin{solution}{exo:b2:plant-flowering:7}
$n = \ln(1/0.15)/0.05 = 1.90/0.05 = 38$ 个寒冷日。各段共计
$47$ 天：植株在第三段的第 11 天完成春化。
若 $\alpha = 0.02$，则需 95 天，植株不会完成春化。
\end{solution}

\begin{solution}{exo:b2:plant-flowering:8}
阻遏基因被染色质标记沉默，这些标记在每次有丝分裂时被复制，
因此茎的每个细胞都“记得”冬季；在生殖细胞系和早期胚胎中，
标记被擦除，该基因重新表达。这种安排是有益的，因为夏季散落的\omterm{def:b2:angiosperm-reproduction:seed}{种子}
不应继承亲本的开花许可，
而必须等待属于自己的
冬季。
\end{solution}

\begin{solution}{exo:b2:plant-flowering:9}
R：萌发（\omterm{prop:b2:plant-flowering:photoperiod}{光敏色素}转变为远红光吸收形式，
即活性形式 Pfr）。R 随后接 FR：不萌发（又转回 Pr）。
R、FR、R：萌发。R、FR、R、FR：不萌发。只有最后一次闪光起作用，
因为每次转化都是完全的，\omterm{def:b2:angiosperm-reproduction:seed}{种子}读取的是色素的
最终状态。
\end{solution}

\begin{solution}{exo:b2:plant-flowering:10}
$p = 0.7$：$r = 1.28$、$1.42$、$1.40$、$-\infty$，对应 $g = 0.4, 0.6, 0.8,
1$：最优值接近 $g = 0.7$。$p = 0.95$：$1.99$、$2.33$、$2.55$、
$-\infty$：最优值向 $g = 1$ 移动（约 $0.948$），
仅是象征性的储备。沙漠中好年景罕见，有利于重度\omterm{def:b2:angiosperm-reproduction:seed}{休眠}，
并形成大量\omterm{def:b2:angiosperm-reproduction:seed}{种子}库；草甸中多数年份为好年景，
有利于几乎所有\omterm{def:b2:angiosperm-reproduction:seed}{种子}萌发。
\end{solution}

\begin{solution}{exo:b2:plant-flowering:11}
各轮均有 B：花瓣、花瓣、\omterm{def:b2:angiosperm-reproduction:flower}{雄蕊}、\omterm{def:b2:angiosperm-reproduction:flower}{雄蕊}。各轮均有 C（C
抑制 A）：\omterm{def:b2:angiosperm-reproduction:flower}{心皮}、\omterm{def:b2:angiosperm-reproduction:flower}{雄蕊}、\omterm{def:b2:angiosperm-reproduction:flower}{雄蕊}、\omterm{def:b2:angiosperm-reproduction:flower}{心皮}。无 A 且无 B：仅 C 起作用，
各处均为\omterm{def:b2:angiosperm-reproduction:flower}{心皮}——四轮皆为\omterm{def:b2:angiosperm-reproduction:flower}{心皮}。C 还能终止\omterm{def:b2:plant-meristems:meristem}{分生组织}的生长；
若无 C，中心将保持无限生长，并在花内
再形成花，即花中套花。
\end{solution}

\begin{solution}{exo:b2:plant-flowering:12}
日长通过光照与生物钟的重合在多个夜晚间
被累积，诱导需要若干这样的夜晚；低温以累积染色质标记的
方式被累积；\omterm{def:b2:angiosperm-reproduction:seed}{种子}库则通过萌发的分散在多年间累积。
积分器会忽略单个异常的日子、一次一月的解冻，或一个糟糕的年份；
若以瞬时数值设定阈值，
植物会在二月的一个温暖星期里开花，
或\omterm{def:b2:angiosperm-reproduction:seed}{种子}会在第一场雨中全部萌发。
\end{solution}

\begin{solution}{pb:b2:plant-flowering:1}
\textbf{1.} 8 小时的黑夜短于临界值 9.5 小时：
不开花。要在八月出售，种植者必须
用黑布人为延长夜晚。
\textbf{2.} 12 小时的黑夜。8 月 15 日开花意味着诱导
约在 6 月 20 日完成，提前 8 周；十二个诱导性
夜晚须从 6 月 8 日左右开始。
\textbf{3.} 在每个夜晚的中段为作物照明几分钟：
一个 14.5 小时的黑夜变为两段约 7 小时的黑夜，两段均低于
临界长度。闪光即可，因为\omterm{prop:b2:plant-flowering:photoperiod}{光敏色素}在数分钟内即被转化，
植物测量的只是
不间断黑暗的长度。
\textbf{4.} 21:00 的一次闪光留下的黑暗期
为 21:00 至 07:30，即 10.5 小时——长于 9.5 小时：
该段可诱导开花。断点必须使每一段都不超过
临界长度。
\textbf{5.} 缺一个夜晚即打断连续十二个诱导性夜晚的序列：
计数重新开始，开花将推迟约等于已累积天数的时间。
苍耳只需一个诱导性夜晚，
因此早已定型。
\textbf{6.} \omterm{prop:b2:plant-flowering:photoperiod}{光敏色素}不吸收绿光，因此植物看不到这次闪光；
红光使其转化为活性形式，被读作白昼；随后的远红光
又将其转回，
因此黑夜未被打断。
\textbf{7.} \omterm{prop:b2:plant-flowering:photoperiod}{长日植物}开花：\omterm{def:b2:plant-flowering:florigen}{成花素}经韧皮部穿过嫁接
愈合处。这说明存在一种可移动、可通过嫁接传递的
信号，为两类日长植物所共有。
\textbf{8.} $n = \ln 10/0.06 = 38.4$：需 39 个寒冷日。
\textbf{9.} 十一月为 10 个寒冷日，一月底累计 35 个，
第 39 个寒冷日是二月的第四个寒冷日：
植株于二月初完成春化。
\textbf{10.} 六天：$f = e^{-0.36} = 0.70 > 0.10$：植株
整个夏季保持营养生长，不结实——冬小麦须
于秋季播种。
\textbf{11.} 若无阻遏物，植株无需低温即可开花：可于春季
播种，同一夏季收获——即春小麦。
\textbf{12.} 阻遏基因染色质上的沉默标记，
在每次有丝分裂时被复制，因此茎尖在整个春季都“记得”冬季；
在\omterm{def:b2:meiosis-heredity:meiosis}{减数分裂}和胚胎中这些标记被擦除，每一代
都须计算属于自己的冬季。
\textbf{13.} 这一记忆必须存在于将来形成花、
并能熬过冬季的细胞中——即\omterm{def:b2:plant-meristems:meristem}{分生组织}——它是这些细胞
染色质的一种状态，而非可移动的信号；测量日长的叶片
会脱落，转而输出一种蛋白质。
\textbf{14.} C 类：萼片、花瓣、花瓣、萼片（A 扩展到第四轮，
花在内部重复出现）。
\textbf{15.} 无\omterm{def:b2:angiosperm-reproduction:flower}{雄蕊}，故无\omterm{def:b2:plant-life-cycles:heterospory}{花粉}；在 C 活性较弱处形成少量\omterm{def:b2:angiosperm-reproduction:flower}{心皮}，
因此偶尔结实。
\textbf{16.} A 类：\omterm{def:b2:angiosperm-reproduction:flower}{心皮}、\omterm{def:b2:angiosperm-reproduction:flower}{雄蕊}、\omterm{def:b2:angiosperm-reproduction:flower}{雄蕊}、\omterm{def:b2:angiosperm-reproduction:flower}{心皮}。
\textbf{17.} C 抑制 A：第一轮成为\omterm{def:b2:angiosperm-reproduction:flower}{心皮}，第二轮（B 与 C）为\omterm{def:b2:angiosperm-reproduction:flower}{雄蕊}——
\omterm{def:b2:angiosperm-reproduction:flower}{心皮}、\omterm{def:b2:angiosperm-reproduction:flower}{雄蕊}、\omterm{def:b2:angiosperm-reproduction:flower}{雄蕊}、\omterm{def:b2:angiosperm-reproduction:flower}{心皮}，即 A \omterm{def:b2:genome-diversification:mutation}{突变}体的花。
\textbf{18.} ABC 基因是一个\omterm{prop:b2:cell-differentiation:transcriptional}{转录因子}家族，由所有开花植物的
共同祖先遗传而来；该模块在蔷薇、金鱼草和 \emph{Arabidopsis}
\omterm{def:b2:cell-differentiation:differentiation}{分化}之前即已存在，
各类群均保留了它。
\textbf{19.} 只有 B 而无 A 或 C，不能决定任何花器官：
形成类似叶片或苞片的器官。
\textbf{20.} $p = 0.8$：$r = 1.44$（$g = 0.5$）、$1.59$（$0.7$）、$1.55$
（$0.9$）。
\textbf{21.} 三者中，$g = 0.7$ 为最佳（真正的最优值接近
$0.8$）；当 $g = 1$ 时，坏年份因子为零，该谱系将在第一个
坏年份中灭绝。
\textbf{22.} $p = 0.5$：$0.51$、$0.41$、$-0.03$：$g = 0.5$ 最佳，
而更低的值可能更好——沙漠中应有更多\omterm{def:b2:angiosperm-reproduction:seed}{休眠}。
\textbf{23.} 种植者的\omterm{def:b2:angiosperm-reproduction:seed}{种子}是贮存而非埋藏的，且确定能够存活；
他无须在田间分散风险，但须在谷仓中保留储备，以便在
歉收的年份补种——这与土壤中的\omterm{def:b2:angiosperm-reproduction:seed}{种子}库是同一种储备，
只是由他而非土壤保管。
\textbf{24.} 小\omterm{def:b2:angiosperm-reproduction:seed}{种子}所含的储备只够一株长一两厘米的幼苗；
若在土壤下的黑暗中萌发，它会在到达光线之前死去，
故须等待光照——而光线只在近地表处才能到达它。
翻耕会把埋藏的\omterm{def:b2:angiosperm-reproduction:seed}{种子}翻上来，
引发杂草的集中萌发。
\textbf{25.} 临界夜长 9.5 小时，覆盖自约 6 月 8 日起至
8 月 15 日开花的时段；39 个寒冷日；最优 $g \approx 0.8$（$p = 0.8$ 时），
约为 $0.5$（$p = 0.5$ 时）。
\end{solution}
